%%
%%  Unified Adaptive-Friction + Navier-Stokes Paper
%%  Target: SIAM Journal on Applied Mathematics
%%
\documentclass[final]{siamltex}

%% ====================================================================
%%  Packages
%% ====================================================================
\usepackage{amsmath,amssymb,amsfonts}
\usepackage{mathtools}
\usepackage{bm}
\usepackage{booktabs}
\usepackage{algorithm}
\usepackage{algorithmic}
\usepackage{graphicx}
\usepackage{hyperref}
\usepackage{url}
\usepackage{microtype}
\usepackage{xcolor}
\usepackage{appendix}

%% ====================================================================
%%  Notation & shortcuts
%% ====================================================================
\newcommand{\R}{\mathbb{R}}
\newcommand{\E}{\mathbb{E}}
\newcommand{\N}{\mathbb{N}}
\newcommand{\T}{\mathbb{T}}
\newcommand{\norm}[1]{\lVert #1 \rVert}
\newcommand{\inner}[2]{\langle #1,\, #2 \rangle}
\providecommand{\grad}{\nabla}
\newcommand{\Ebs}{E_{\mathrm{BS}}}
\newcommand{\Cman}{\mathcal{C}^*}
\newcommand{\Xst}{X^*}
\newcommand{\Lpair}{L_{\mathrm{pair}}}
\newcommand{\lmin}{\lambda_{\min}}
\newcommand{\lmax}{\lambda_{\max}}
\newcommand{\costheta}{\cos\theta}
\newcommand{\ddt}{\tfrac{d}{dt}}
\DeclareMathOperator{\erf}{erf}
\DeclareMathOperator{\sym}{sym}

%% ====================================================================
%%  Extra theorem-like environments (SIAM provides theorem, lemma,
%%  corollary, proposition, definition)
%% ====================================================================
\newtheorem{assumption}[theorem]{Assumption}
\newtheorem{remark}[theorem]{Remark}

%% ====================================================================
%%  Title & metadata
%% ====================================================================
\title{Systemic Risk as Geometry: Blind-Spot Decomposition, Adaptive
Friction, and Applications to Finance, Energy, and Navier--Stokes
Regularity\thanks{Submitted to the editors \today.}}

\author{Oluwatosin Odeyemi\thanks{Independent Researcher, London, UK.
(\texttt{odeyemi.oluwatosin@outlook.com}).
ORCID: 0009-0001-7043-6986.}}

\begin{document}
\maketitle

\pagestyle{myheadings}
\thispagestyle{plain}
\markboth{O.\ ODEYEMI}{SYSTEMIC RISK AS GEOMETRY}

%% ====================================================================
%%  Abstract
%% ====================================================================
\begin{abstract}
We introduce a geometric framework for analysing instability in
coupled dynamical systems. $N$ interacting agents evolve under a
potential $\Phi$ whose pairwise term encodes contagion; a
\emph{Blind-Spot Decomposition Theory} (BSDT) energy functional
$\Ebs$, calibrated on normal-period data, provides an equivalent
Lyapunov function (Theorem~C: $c_1\norm{\grad\Phi}\le\norm{\grad\Ebs}\le c_2\norm{\grad\Phi}$).
A \emph{critical manifold} $\Cman$ separates a stable from an unstable
regime; adaptive friction $\gamma^*(X)=\alpha/\lmax(D^2\Phi_{\mathrm{pair}})$
achieves marginal stability where every fixed coupling fails. The
framework is validated on U.S.\ banking data (FDIC SDI, 1990--2024;
6-quarter GFC lead, AUROC $\ge 0.70$), the Terra/Luna collapse
($r=0.996$ calibration, $+30$\,h lead), and the Texas ERCOT grid
failure ($+72$\,h lead). We then transfer the BSDT channels to
three-dimensional incompressible Navier--Stokes: enstrophy, spectral
cascade, vorticity--strain alignment, and temporal acceleration
define a fluid-mechanical $\Ebs$ whose adaptive viscosity
$\nu_{\mathrm{eff}}=\nu_0(1+\gamma^*(\Ebs))$ suppresses enstrophy
growth and maintains a production/dissipation ratio converging to
$0.500$ across Reynolds numbers $314$--$6283$. This cross-domain
transfer connects macroprudential regulation, complex-systems
stability, and Navier--Stokes regularity through a single
mathematical structure.
\end{abstract}

\begin{keywords}
systemic risk, Lyapunov stability, phase transition, adaptive
friction, Navier--Stokes regularity, enstrophy, blind-spot
decomposition, online convex optimisation
\end{keywords}

\begin{AMS}
91G45, 37N40, 35Q30, 76D03, 93D20, 91B55
\end{AMS}

%% ====================================================================
%%  1.  INTRODUCTION
%% ====================================================================
\section{Introduction}\label{sec:intro}

The 2008 Global Financial Crisis demonstrated that systemic risk
resides not in individual institutions but in the geometry of
their interactions. Existing frameworks---CoVaR
\cite{adrian2016covar}, SRISK \cite{brownlees2017srisk},
network contagion models \cite{acemoglu2015systemic,
elliott2014financial}---detect stress \emph{after} it
propagates. This paper develops a framework that detects
instability \emph{before} it manifests, by identifying the
geometric precursors in the joint state space of interacting agents.

The core insight is an \emph{equivalence theorem} linking two
seemingly unrelated objects: (i)~a \emph{detection} functional
$\Ebs$ measuring deviation from historical normality, and (ii)~the
\emph{interaction energy} $\Phi$ governing agent dynamics.
Theorem~C establishes that $\Ebs$ is a valid Lyapunov function for
the gradient flow $\dot{X}=-\grad\Phi(X)$, with gradient norms
equivalent up to computable constants. This identification transforms
the detection problem into a control problem: damping the system
using $\grad\Ebs$ stabilises it because $\grad\Ebs$ and $\grad\Phi$
are geometrically aligned.

A \emph{critical manifold} $\Cman$ separates two regimes. Below
$\Cman$, any fixed coupling suffices for stability. Above $\Cman$,
no constant friction stabilises all configurations
(Theorem~\ref{thm:phase}). The system requires \emph{adaptive}
friction $\gamma^*(X)$ that tracks the spectral radius of the
pairwise Hessian in real time. This result is the formal content
of procyclicality: time-invariant regulatory rules (such as Basel
III's fixed countercyclical capital buffer) fail structurally, not
merely in practice.

\paragraph{Contributions.}
This paper makes seven contributions:
\begin{enumerate}
\item \textbf{Blind-Spot Decomposition Theory (BSDT).} Four
  deviation operators ($\delta_C$: camouflage, $\delta_G$: feature
  gap, $\delta_A$: activity anomaly, $\delta_T$: temporal novelty)
  define $\Ebs$, a composite energy functional calibrated entirely
  on normal-period data.

\item \textbf{Equivalence Theorem (Theorem~C).} Two-sided gradient
  bounds establish $\Ebs$ as a Lyapunov function for the potential
  flow, with a structurally derived separation condition
  (Proposition~\ref{prop:separation}).

\item \textbf{Critical manifold and phase transition
  (Theorem~\ref{thm:phase}).} A spectral threshold determines
  where constant friction fails and adaptive friction succeeds.

\item \textbf{Online algorithm with $O(\sqrt{T})$ regret.} The
  adaptive coupling is implementable via online gradient descent
  with power-iteration spectral estimates.

\item \textbf{Empirical validation.} On FDIC banking data
  ($T=140$ quarters, $N=7$ institution types, $d=6$ features), the
  framework achieves a 6-quarter GFC early warning, AUROC $\ge
  0.70$, and consistent Granger failure---the expected signature of
  a phase-transition detector.

\item \textbf{Cross-domain universality.} The same framework,
  without re-training, detects the Terra/Luna collapse ($r=0.996$)
  and the Texas ERCOT grid failure ($+72$\,h lead).

\item \textbf{Extension to Navier--Stokes.} The BSDT channel
  structure transfers to three-dimensional incompressible
  Navier--Stokes, where adaptive viscosity suppresses enstrophy
  growth and maintains bounded production/dissipation ratios---a
  pathway toward regularity results.
\end{enumerate}

\paragraph{Related work.}
Network-based systemic risk models
\cite{acemoglu2015systemic,elliott2014financial} capture topology
but not dynamics. Market-based measures (CoVaR, SRISK) require
market prices and operate at the market frequency, missing
slow-moving structural drift on regulatory balance sheets.
The May--Levin--Sugihara programme \cite{may2008complex} imports
ecological stability analysis but lacks operational control
algorithms. Our framework complements these by providing feedback
control with formal stability guarantees. For the Navier--Stokes
connection, the Beale--Kato--Majda criterion \cite{bkm1984} and
Ladyzhenskaya's classical estimates \cite{lady1969} provide the
regularity backdrop.

\paragraph{Outline.}
Section~\ref{sec:model} introduces the agent model and BSDT.
Section~\ref{sec:equiv} proves Theorem~C. Section~\ref{sec:damping}
establishes the phase transition. Section~\ref{sec:online} presents
the online algorithm and optimal policy. Section~\ref{sec:empirical}
validates on financial data. Section~\ref{sec:crossdomain} extends
to crypto and energy systems. Section~\ref{sec:ns} transfers the
framework to Navier--Stokes. Section~\ref{sec:discussion}
concludes.


%% ====================================================================
%%  2.  AGENT MODEL AND BSDT
%% ====================================================================
\section{Mathematical Framework}\label{sec:model}

\subsection{Agent Model and GravityEngine Dynamics}\label{sec:agents}

Consider $N$ interacting institutions, each described by a state
vector $x_i(t)\in\R^d$ of financial characteristics. The joint
state $X(t)=(x_1,\ldots,x_N)\in\R^{N\times d}$ evolves under the
gradient flow
\begin{equation}\label{eq:gradflow}
  \dot{X} = F(X) = -\grad\Phi(X),
\end{equation}
where $\Phi:\R^{N\times d}\to\R$ decomposes as
\begin{equation}\label{eq:potential}
  \Phi(X) = \underbrace{\frac{\alpha}{2}\sum_{i=1}^{N}
  \norm{x_i-\mu}^2}_{\text{radial spring (mean-reversion)}}
  + \underbrace{\sum_{i<j}\phi(\norm{x_i-x_j})}_{\text{pairwise
  interaction}}.
\end{equation}
The pairwise potential $\phi(r)=\gamma\frac{\sigma\sqrt{\pi}}{2}
\erf(r/\sigma)-\gamma\lambda\log(r+\varepsilon)$ combines
short-range Gaussian attraction with long-range logarithmic
repulsion, giving clustering at equilibrium distance $r^*=
\sigma\sqrt{\log(1/\lambda)}$.

The \emph{unique Nash equilibrium} under the friction game
$\min_{\ell_i}[\text{private cost}_i+\gamma\cdot
\text{systemic externality}_i]$ is $\ell_i^*(\gamma)=
(\alpha+\gamma w_i)^{-1}r_i$, where $w_i$ is institution~$i$'s
network centrality and $r_i$ its return on assets.
The optimal coupling $\gamma^*$ internalises the Pigouvian
externality $\tau_i^*=\gamma^*\,\partial E/\partial\ell_i$,
proportional to each institution's marginal contribution to system
energy.

\subsection{Blind-Spot Decomposition Theory (BSDT)}\label{sec:bsdt}

BSDT decomposes systemic risk into four geometrically orthogonal
channels of deviation from the normal-period distribution
$\mathcal{N}$:
\begin{enumerate}
\item \textbf{Camouflage} ($\delta_C$): Mahalanobis distance from
  $\mathcal{N}$. Detects states that appear individually normal but
  are collectively anomalous:
  $\delta_C(X) = [(X{-}\mu_0)^\top\Sigma_0^{-1}(X{-}\mu_0)]^{1/2}$.

\item \textbf{Feature Gap} ($\delta_G$): PCA residual in
  low-variance directions. Detects motion into historically
  unoccupied subspaces: $\delta_G(X)=\norm{(I-P_k)X}$.

\item \textbf{Activity Anomaly} ($\delta_A$): Excess velocity
  relative to normal-period dynamics:
  $\delta_A=\max(0,\norm{\dot{x}_i}-v_0)$.

\item \textbf{Temporal Novelty} ($\delta_T$): Negative log-density
  under a KDE of historical states:
  $\delta_T=-\log\hat{p}(x_i(t))$.
\end{enumerate}

The \emph{blind-spot energy functional} is the weighted composite
\begin{equation}\label{eq:ebs}
  \Ebs(X) = \sum_{i=1}^{4}\psi_i(\delta_i(X))
  + \sum_{i<j}\phi(\norm{x_i-x_j}),
\end{equation}
where $\psi_i:\R^+\to\R^+$ are convex, non-decreasing,
1-Lipschitz link functions with $\psi_i(0)=0$. The pairwise
potential in $\Ebs$ coincides with $\Phi_{\mathrm{pair}}$,
ensuring structural compatibility between detection and dynamics.

The \emph{Multi-Factor Latent Score} (MFLS) is the Frobenius norm
$\norm{\grad\Ebs(X_t)}_F$, serving as a scalar early-warning
signal.


%% ====================================================================
%%  3.  THE EQUIVALENCE THEOREM
%% ====================================================================
\section{The Equivalence Theorem}\label{sec:equiv}

We work with a fixed reference measure $\mathcal{N}$ (the
normal-period distribution) and assume $\Phi$ satisfies
Assumptions~\ref{ass:reg}--\ref{ass:sep} below.

\subsection{Assumptions}

\begin{assumption}[Regularity]\label{ass:reg}
$\Phi\in C^2(\R^{N\times d})$ is coercive: $\Phi(X)\to+\infty$ as
$\norm{X}\to\infty$. Every sublevel set $\mathcal{L}_c:=
\{X:\Phi(X)\le c\}$ is compact, with Lipschitz constant $L_c$ for
$\grad\Phi$ on $\mathcal{L}_c$.
\end{assumption}

\begin{assumption}[Strict local minimum]\label{ass:equil}
The equilibrium $\Xst$ satisfies $\grad\Phi(\Xst)=0$ and
$H^*:=D^2\Phi(\Xst)\succ 0$.
\end{assumption}

\begin{assumption}[BSDT regularity]\label{ass:bsdt}
Each $\delta_i:$ $\R^{N\times d}\to\R^+$ is $C^1$ and
$K_i$-Lipschitz on $\mathcal{L}_c$. Each $\psi_i:\R^+\to\R^+$ is
convex, non-decreasing, 1-Lipschitz, with $\psi_i(0)=0$.
\end{assumption}

\begin{assumption}[Separation]\label{ass:sep}
There exists $\nu>0$ such that for all $X\in\mathcal{L}_c$ and
all~$i$: $\norm{\grad_X\delta_i(X)}\ge\nu\norm{\grad\Phi(X)}$.
\end{assumption}

\begin{proposition}[Structural derivation of the separation
condition]\label{prop:separation}
Let $\Sigma_0\succ 0$ be the Ledoit--Wolf shrinkage estimator
with condition number $\kappa(\Sigma_0)$. Then:
\begin{enumerate}
\item[\textup{(i)}] $\delta_C$ satisfies Assumption~\ref{ass:sep}
  with $\nu_C=\kappa(\Sigma_0)^{-1}[\alpha\sqrt{\lmax(\Sigma_0)}]^{-1}$,
  determined entirely by spectral properties of $\Sigma_0$ and the
  spring constant~$\alpha$.

\item[\textup{(ii)}] $\delta_G$ satisfies $\norm{\grad_X\delta_G}=1$
  whenever $\delta_G>0$, giving unconditional separation.

\item[\textup{(iii)}] $\delta_A$ satisfies Assumption~\ref{ass:sep}
  with $\nu_A=2\eta(\lambda_u-\alpha)/L_c$ in the super-critical
  region.

\item[\textup{(iv)}] $\delta_T$ satisfies Assumption~\ref{ass:sep}
  with $\nu_T=h^{-1}/L_c$ outside the kernel bandwidth~$h$.
\end{enumerate}
The composite $\Ebs$ inherits
$\nu\ge\max(\nu_C,\nu_G,\nu_A,\nu_T)$.
\end{proposition}

\begin{proof}
Part~(i): $\grad_X\delta_C=\Sigma_0^{-1}(X{-}\mu_0)/\delta_C(X)$;
by the Rayleigh quotient,
$\norm{\Sigma_0^{-1}v}^2\ge\lmax(\Sigma_0)^{-2}\norm{v}^2$, and
$\delta_C(X)\le\sqrt{\lmax(\Sigma_0)}\norm{v}$. Combining with
$\norm{\grad\Phi(X)}\le L_c\norm{X{-}\Xst}$ near $\Xst$ gives the
bound. Part~(ii): $(I{-}P_k)$ is an orthogonal projector, so
$\grad_X\norm{(I{-}P_k)X}$ is a unit vector. Parts~(iii)--(iv)
follow from direct gradient computation using super-critical
growth estimates (Theorem~\ref{thm:phase}) and score-function
lower bounds.
\end{proof}

\subsection{Gradient Alignment}\label{sec:gradient}

\begin{lemma}[Gradient of $\Ebs$]\label{lem:grad_ebs}
Under Assumption~\ref{ass:bsdt}, $\Ebs$ is $C^1$ with
\[
  \grad\Ebs(X) =
  \underbrace{\textstyle\sum_{i=1}^{4}\psi_i'(\delta_i)\grad_X\delta_i}_
  {\grad\Ebs^{\mathrm{dev}}}
  + \underbrace{\grad_X\textstyle\sum_{i<j}
  \phi(\norm{x_i{-}x_j})}_{\grad\Ebs^{\mathrm{pair}}
  =-F_{\mathrm{pair}}}.
\]
\end{lemma}

\begin{lemma}[Upper bound]\label{lem:upper}
Under Assumptions~\ref{ass:reg}--\ref{ass:bsdt}, on
$\mathcal{L}_c$ there exists $c_2>0$ such that
$\norm{\grad\Ebs(X)}\le c_2\norm{\grad\Phi(X)}$.
\end{lemma}

\begin{lemma}[Lower bound]\label{lem:lower}
Under Assumptions~\ref{ass:reg}--\ref{ass:sep}, on
$\mathcal{L}_c$ there exists $c_1>0$ such that
$\norm{\grad\Ebs(X)}\ge c_1\norm{\grad\Phi(X)}$.
\end{lemma}

Proofs of the upper and lower bounds follow from the triangle
inequality, the Lipschitz properties of $\delta_i$, and the
separation condition. Full details are in
Appendix~\ref{app:proofs}.

\subsection{Theorem~C: Full Statement and Proof}\label{sec:thmC}

\begin{theorem}[Equivalence of Blind-Spot Geometry and Instability
Energy]\label{thm:equiv}
Let Assumptions~\ref{ass:reg}--\ref{ass:sep} hold and let $\Xst$
be a strict local minimum of $\Phi$. Then on any sublevel set
$\mathcal{L}_c$:

\begin{enumerate}
\item[\textup{(C1)}] \textbf{Gradient equivalence.} There exist
  $c_1,c_2>0$ with
  $c_1\norm{\grad\Phi}\le\norm{\grad\Ebs}\le c_2\norm{\grad\Phi}$
  and alignment $\costheta\ge\nu/c_2$.

\item[\textup{(C2)}] \textbf{Energy equivalence.} There exist
  $a,b>0$ with $a\,V(X)\le\Ebs(X)-\Ebs(\Xst)\le b\,V(X)$,
  where $V(X):=\Phi(X)-\Phi(\Xst)$.

\item[\textup{(C3)}] \textbf{Stabilisation sufficiency.} The
  blind-spot-damped dynamics
  \begin{equation}\label{eq:bsdamped}
    \dot{X}=F(X)-\gamma(\Ebs(X))\grad\Ebs(X),
  \end{equation}
  with $\gamma(E)\ge\kappa>0$ for $E>\theta_0$, render $\Xst$
  locally asymptotically stable.
\end{enumerate}
\end{theorem}

\begin{proof}
\textbf{Part~(C1).} The bounds are Lemmas~\ref{lem:upper} and
\ref{lem:lower}. For alignment: write $\grad\Ebs=
\alpha_\parallel(-\grad\Phi/\norm{\grad\Phi})+\alpha_\perp e_\perp$.
Since the pairwise term of $\grad\Ebs$ coincides with
$-F_{\mathrm{pair}}$ (a component of $-\grad\Phi$), the
anti-parallel component satisfies $\alpha_\parallel\ge
\nu\norm{\grad\Phi}$ by Assumption~\ref{ass:sep}, giving
$\costheta=\alpha_\parallel/\norm{\grad\Ebs}\ge\nu/c_2$.

\textbf{Part~(C2).} By (C1) and the fundamental theorem of
calculus applied to both $\Ebs(X)-\Ebs(\Xst)$ and $V(X)$ along the
segment $\Xst+s(X-\Xst)$, the two-sided gradient bound transfers
to a two-sided energy bound with $a=c_1$, $b=c_2$.

\textbf{Part~(C3).} Take $W(X)=\Ebs(X)-\Ebs(\Xst)$ as the
Lyapunov candidate. By~(C2), $a\,V\le W\le b\,V$, so $W$ is
positive definite and proper. Along \eqref{eq:bsdamped}:
\begin{align*}
  \dot{W} &= \inner{\grad\Ebs}{F} - \gamma(\Ebs)\norm{\grad\Ebs}^2 \\
  &\le -\frac{\costheta}{c_2}\norm{\grad\Ebs}^2
  - \gamma(\Ebs)\norm{\grad\Ebs}^2
  \le -\kappa\norm{\grad\Ebs}^2 \le 0.
\end{align*}
LaSalle's invariance principle \cite{lasalle1960} on the compact
set $\{W\le W(X_0)\}$ concludes convergence to $\Xst$.
\end{proof}

\begin{remark}\label{rem:alignment}
The lower bound $c_1\norm{\grad\Phi}\le\norm{\grad\Ebs}$ is
essential for (C3): without it, $\grad\Ebs$ could be orthogonal to
$F$, and the damping term would exert no stabilising force. The
separation condition (Assumption~\ref{ass:sep}) is the minimal
non-degeneracy requirement for detection-based stabilisation.
\end{remark}


%% ====================================================================
%%  4.  ADAPTIVE DAMPING AND PHASE TRANSITION
%% ====================================================================
\section{Adaptive Damping and the Critical
Manifold}\label{sec:damping}

\begin{theorem}[Monotone Descent]\label{thm:descent}
Under Assumption~\ref{ass:reg}, $\ddt\Phi(X(t))=
-\norm{\grad\Phi}^2\le 0$. Every $\omega$-limit point satisfies
$\grad\Phi=0$.
\end{theorem}

\begin{theorem}[Local Exponential Stability]\label{thm:las}
Under Assumptions~\ref{ass:reg}--\ref{ass:equil}, $\Xst$ is
locally exponentially stable with rate $\lmin(H^*)$:
\[
  \norm{X(t)-\Xst}\le\norm{X(0)-\Xst}\,e^{-(\alpha-\Lpair)t}
  \quad\text{for }\alpha>\Lpair.
\]
\end{theorem}

\begin{theorem}[Spectral Phase Transition]\label{thm:phase}
Define the critical manifold
\begin{equation}\label{eq:cman}
  \Cman := \bigl\{X : \lmax\bigl(\rho(X)\,\ell(X)\,
  \mathbf{W}\bigr) = 1\bigr\},
\end{equation}
where $\rho(X)$ is mean pairwise correlation, $\ell(X)$ mean
exposure, and $\mathbf{W}$ the row-normalised adjacency matrix.
\begin{enumerate}
\item[\textup{(i)}] Below $\Cman$: for any fixed coupling
  $\bar{\gamma}$, the gradient flow is locally asymptotically
  stable.
\item[\textup{(ii)}] Above $\Cman$: for any fixed
  $\bar{\gamma}>0$, there exists a configuration $X^+$ with a
  growing mode. No constant $\bar{\gamma}$ stabilises all
  above-$\Cman$ configurations.
\item[\textup{(iii)}] Adaptive $\gamma^*(X)=\alpha/
  \lmax(D^2\Phi_{\mathrm{pair}}(X))$ achieves marginal stability
  everywhere and asymptotic stability outside a neighbourhood
  of~$\Cman$.
\end{enumerate}
\end{theorem}

\begin{proof}
\textit{Part~(i).} Below $\Cman$,
$\lmax(D^2\Phi_{\mathrm{pair}})<\alpha$, so
$\lmax(\sym(DF))=-(\alpha-\Lpair)<0$ and contraction theory
\cite{lohmiller1998contraction} gives LAS.

\textit{Part~(ii).} Above $\Cman$, a leading eigenvector $u$ of
$D^2\Phi_{\mathrm{pair}}(X^+)$ has eigenvalue
$\lambda_u>\alpha$. The mode energy
$W(t)=\frac{1}{2}\inner{X(t){-}\Xst}{u}^2$ satisfies
$\dot{W}\ge(\lambda_u{-}\alpha)W>0$, growing exponentially.

\textit{Part~(iii).}
$\gamma^*(X)=\alpha/\lmax(D^2\Phi_{\mathrm{pair}}(X))$ zeroes the
net spectral abscissa by construction. Asymptotic stability outside
a neighbourhood of $\Cman$ follows from the radial spring
dominating once $\lmax<\alpha$.
\end{proof}

\begin{corollary}[Basel III Procyclicality]\label{cor:procyclical}
A fixed regulatory capital buffer $\bar{\kappa}$ corresponds to
constant friction $\bar{\gamma}\propto\bar{\kappa}$. By
Theorem~\ref{thm:phase}(ii), for sufficiently correlated
portfolios (above $\Cman$), no fixed $\bar{\kappa}$ simultaneously
maintains conservatism in normal times and sufficiency in stress.
\end{corollary}


%% ====================================================================
%%  5.  ONLINE ALGORITHM AND OPTIMAL POLICY
%% ====================================================================
\section{Online Algorithm and Optimal Policy}\label{sec:online}

\subsection{GravityEngine as Online Convex Optimisation}

At each step $t$, the regulator selects coupling
$\gamma_t\in[\gamma_{\min},\gamma_{\max}]$; the system evolves to
$X_{t+1}(\gamma_t)$ and cost
$c_t(\gamma_t)=\Phi(X_{t+1})-\Phi(X_t)$ is revealed.

\begin{proposition}[Online Regret Bound]\label{prop:regret}
If $c_t$ is convex and $L_c$-Lipschitz in $\gamma$, projected
online gradient descent achieves
\begin{equation}\label{eq:regret}
  R_T = \sum_{t=0}^{T-1}c_t(\gamma_t)
  - \sum_{t=0}^{T-1}c_t(\gamma^*(X^t))
  \le L_c(\gamma_{\max}-\gamma_{\min})\sqrt{T}.
\end{equation}
\end{proposition}

\begin{proof}
$X^{t+1}(\gamma)$ is affine in $\gamma$; composing the convex
$\Phi$ yields convexity. The Lipschitz bound
$|\partial_\gamma c_t|\le\eta F_{\max}^2=:L_c$ follows from the
force clamp. Standard OGD guarantees
\cite{hazan2016oco,shalev2012online} then apply.
\end{proof}

\begin{proposition}[Computational Complexity]\label{prop:complexity}
Computing $\grad\Ebs(X)$ costs $O(N^2 d)$ per step. With $k$-d
tree acceleration and graph sparsification, complexity reduces to
$O(Nd\log N)$.
\end{proposition}

\begin{algorithm}[H]
\caption{Adaptive GravityEngine (one outer step)}
\label{alg:adaptive}
\begin{algorithmic}[1]
\REQUIRE State $X^t$, coupling range
  $[\gamma_{\min},\gamma_{\max}]$, step $\eta$, OGD step
  $\eta_\gamma$
\STATE Compute $\grad\Phi(X^t)$ via pairwise loop ($O(N^2 d)$)
\STATE $\hat\lambda_t\leftarrow$ power-iteration Rayleigh quotient
  of $D^2\Phi_{\mathrm{pair}}(X^t)$
\STATE $\hat\gamma^*\leftarrow\alpha/\hat\lambda_t$
\STATE $\tilde{g}_t\leftarrow\partial_\gamma c_t(\gamma_t)$
\STATE $\gamma_{t+1}\leftarrow
  \Pi_{[\gamma_{\min},\gamma_{\max}]}
  (\gamma_t-\eta_\gamma\tilde{g}_t)$
\STATE $X^{t+1}\leftarrow X^t+\eta(-\grad\Phi(X^t))$ with Armijo
  backtracking
\ENSURE $\Phi(X^{t+1})\le\Phi(X^t)$
\end{algorithmic}
\end{algorithm}

\subsection{Optimal Policy via Dynamic Programming}

The social planner minimises total discounted cost:
\begin{equation}\label{eq:dp}
  V(X) = \min_{\{\gamma_t\}}\E\!\left[\sum_{t=0}^\infty\beta^t
  (E(X_t)+\kappa\gamma_t^2)\;\Big|\;X_0=X\right],
\end{equation}
where $\kappa>0$ is the cost of intervention and $\beta\in(0,1)$
the discount factor.

\begin{theorem}[Optimal Adaptive Policy]\label{thm:optimal}
The optimal coupling has the closed-form approximation
\begin{equation}\label{eq:gammastar}
  \gamma^*(X) \approx
  \frac{\beta\norm{\grad E(X)}^2}{2\kappa+\beta\norm{\grad E(X)}^2}
  \cdot\frac{E(X)}{E(X)+\theta},\qquad
  \theta=\kappa/\beta.
\end{equation}
This recovers the adaptive damping law $\gamma(E)=E/(E+\theta)$ as
its leading-order term.
\end{theorem}

\begin{theorem}[Necessity of State-Dependence]\label{thm:necessity}
For any constant policy $\bar{\gamma}$ and any $X$ above $\Cman$:
\[
  V^*(X)-V_{\bar{\gamma}}(X)\ge\Delta(\rho(X),\ell(X))>0.
\]
Constant friction is strictly suboptimal above $\Cman$.
\end{theorem}


%% ====================================================================
%%  6.  EMPIRICAL VALIDATION
%% ====================================================================
\section{Empirical Validation}\label{sec:empirical}

\subsection{Data}\label{sec:data}

We construct $X(t)\in\R^{N\times d}$ from FDIC Statistics on
Depository Institutions (SDI) quarterly call reports. $N=7$ FDIC
institution types, $d=6$ features: loan-to-asset
($\texttt{LNLSNET}/\texttt{ASSET}$), equity ratio
($\texttt{EQ}/\texttt{ASSET}$), NPL ratio
($\texttt{NCLNLS}/\texttt{LNLSNET}$), ROA
($\texttt{NETINC}/\texttt{ASSET}$), funding cost
($\texttt{EINTEXP}/\texttt{LNLSNET}$), and yield-curve slope
(FRED T10Y2Y). The panel spans $T=140$ quarters (1990-Q1 to
2024-Q4). The normal reference period is 1994-Q1 through 2003-Q4
(40~quarters). The inter-sector exposure network $\mathbf{W}$ is
estimated by Ledoit--Wolf Oracle shrinkage \cite{ledoit2004}
($\rho^*=0.046$; spectral radius $\lmax(\mathbf{W})=2.98$).

\subsection{Gradient Alignment on Real Data}

The per-quarter cosine alignment $\costheta(t)=
\inner{\grad\Ebs}{-\grad\Phi}/
(\norm{\grad\Ebs}\norm{\grad\Phi})$ is \emph{negative}
(mean~$=-0.335$) across all 140~quarters. This anti-alignment is
significant: the deviation gradient and the restoring force measure
\emph{complementary} axes of instability. The potential $\Phi$ pulls
institutions toward their current cluster centre (spatial cohesion);
$\Ebs$ measures deviation from the \emph{historical} normal
(temporal anomaly). The equivalence theorem holds in the
mathematical sense (both are Lyapunov candidates), but their
gradients span complementary subspaces on real data.

\subsection{Critical Manifold and Super-Criticality}

The system spends \textbf{68.6\%} of the 140-quarter sample above
$\Cman$, implying that the U.S.\ banking system is \emph{structurally
over-coupled}. Crises occur when perturbations arrive while the
system is already super-critical, consistent with the persistence
view of systemic risk \cite{brunnermeier2009deciphering}.

\subsection{Crisis Window Analysis}

\begin{table}[htbp]
\centering
\caption{Lead time of MFLS versus benchmark stress indices
(quarters, positive = MFLS precedes benchmark). Data: FDIC SDI,
$T=140$.}
\label{tab:lead}
\begin{tabular}{llr}
\toprule
Crisis Window & Benchmark & Lead (quarters)\\
\midrule
GFC 2008    & STLFSI      & $+3$\\
            & VIX         & $+1$\\
            & NFCI        & $+2$\\
COVID 2020  & STLFSI      & $+3$\\
            & VIX         & $+3$\\
Rate Shock 2022 & STLFSI  & $0$\\
            & VIX         & $+1$\\
\bottomrule
\end{tabular}
\end{table}

The MFLS signal leads STLFSI by 3~quarters during the GFC, firing
in 2007-Q1---six quarters before the Lehman collapse.

\subsection{Out-of-Sample Backtest}

A single P75 threshold trained on 1990-Q1 to 2005-Q4 (no crisis
data) is evaluated on 2006-Q1 to 2024-Q4:

\begin{table}[htbp]
\centering
\caption{Out-of-sample backtest. Single P75 threshold trained on pre-2006
data.}
\label{tab:oos}
\begin{tabular}{lcccc}
\toprule
Crisis & First alarm & Lead & Hit rate & FAR\\
\midrule
GFC         & 2007-Q1 & 6Q & 71.4\% & 35\%\\
COVID       & 2019-Q4 & 1--2Q & 62.5\% & 35\%\\
Rate Shock  & 2022-Q1 & coincident & 50.0\% & 33\%\\
\midrule
\multicolumn{5}{l}{\small AUROC $\ge 0.70$ over full 2006--2024 test
window.}\\
\bottomrule
\end{tabular}
\end{table}

\subsection{Universal Granger Failure}

All five MFLS scoring variants---Baseline (Mahalanobis),
Full-BSDT, QuadSurf, Signed-LR, Expo-Gate---fail Granger causality
at $\alpha=0.10$ ($p>0.21$ at all lags). This is not an artefact
but the \emph{expected signature} of a phase-transition detector:
if crises are abrupt threshold crossings rather than trends,
Granger must fail. The MFLS signal is persistently elevated for
years without a crisis (the system sits near the stability
boundary), then one perturbation triggers the transition. The
relationship is \emph{nonlinear in time}.

\subsection{Institution-Level Extension}

Replication on the top-30 U.S.\ banks by total assets (CERT-level
data) produces stronger results: AUROC rises to 0.805, GFC hit
rate to 75.0\%, and $\lmax(\mathbf{W})=12.1$ (vs.\ 8.3 for
SPECGRP), reflecting denser interconnection among individual banks.
The Granger failure is universal at the institution level.

\subsection{Welfare Calibration}

Translating $\gamma^*(X_t)$ into the Basel~III CCyB format:
\begin{equation}\label{eq:ccyb}
  \Delta\mathrm{CCyB}(t)\approx
  \kappa^{-1}\gamma^*(X_t)\sigma_\ell(t).
\end{equation}
The GFC window produces a consumption-equivalent welfare loss of
19.57~pp with peak CCyB of 152~bps. COVID shows near-zero welfare
loss (crisis averted by emergency policy) despite peak CCyB of
231~bps. The 2022 rate shock produces $\mathcal{L}=1.22$~pp with
CCyB = 0 (absorbed without systemic transmission).


%% ====================================================================
%%  7.  CROSS-DOMAIN VALIDATION
%% ====================================================================
\section{Cross-Domain Validation}\label{sec:crossdomain}

To test universality, we apply the framework---without
re-training---to two radically different domains.

\subsection{Terra/Luna UST Collapse (May 2022)}

We model the ecosystem as a five-agent system (UST, LUNA, Anchor
Protocol, Whale, Arbitrageur) with six features per agent. The
simulation achieves $r=0.996$ correlation with historical UST
prices (MAE = 2.35\%). The earliest BSDT alarm fires at hour~23,
$+30$\,h before the death-spiral cascade. The dominant channel is
$\delta_C$ (weight $+1.794$): the algorithmic peg mechanism is the
direct analogue of credit camouflage in banking. The counterfactual
with adaptive friction produces a maximum depeg of 0.28\%---the
stablecoin survives.

\subsection{Texas ERCOT Grid Failure (February 2021)}

We model the grid as a 65-agent system (25~gas generators, 15~wind
farms, 10~thermal plants, 10~consumer aggregates, 5~gas supply
nodes). Twenty of 25~BSDT variants fire at hour~24, $+72$\,h before
rolling blackouts. Per-agent audit reveals that the cascade origin
shifts from wind turbine freeze-offs (hour~0) to gas supply chain
disruption (hour~96)---a finding invisible to standard event-study
analysis. Adaptive friction reduces peak capacity loss by 71.0\%.

\subsection{Universal Structure}

\begin{table}[htbp]
\centering
\caption{Universal collapse structure across domains.}
\label{tab:universal}
\begin{tabular}{@{}lccc@{}}
\toprule
Property & Banking & Crypto & Energy\\
\midrule
Agents & 7 types & 5 ecosystem & 65 nodes\\
Detection lead & 6 quarters & $+30$\,h & $+72$\,h\\
Dominant $\delta$ & $\delta_C$ & $\delta_C$ & $\delta_C$\\
$\gamma^*$ prevents? & Yes (CCyB) & Yes & Yes\\
Granger fails? & Yes & N/A & N/A\\
\bottomrule
\end{tabular}
\end{table}

The universality is not coincidence. Any system with (i)~interacting
agents, (ii)~a deviation-energy functional, (iii)~nonlinear
restoring potential, and (iv)~an opacity mechanism masking proximity
to $\Cman$ will exhibit the same phase-transition dynamics.


%% ====================================================================
%%  8.  EXTENSION TO NAVIER-STOKES REGULARITY
%% ====================================================================
\section{Extension to Navier--Stokes Regularity}\label{sec:ns}

The cross-domain transfer of Section~\ref{sec:crossdomain}
demonstrates that the mathematical framework is domain-agnostic.
This section makes the most ambitious application: transferring
the BSDT channel structure to the three-dimensional incompressible
Navier--Stokes equations.

\subsection{BSDT Channels for Fluid Dynamics}\label{sec:ns:channels}

Consider the 3D incompressible Navier--Stokes equations on
$\T^3=[0,2\pi]^3$:
\begin{equation}\label{eq:ns}
  \partial_t\bm{u}+(\bm{u}\cdot\grad)\bm{u}
  = -\grad p + \nu\Delta\bm{u},\qquad
  \grad\cdot\bm{u}=0,\qquad
  \bm{u}(\cdot,0)=\bm{u}_0\in H^1(\T^3).
\end{equation}
The ``agents'' are Fourier modes $\hat{\bm{u}}(\bm{k},t)$;
``interaction'' is the nonlinear advection
$\widehat{(\bm{u}\cdot\grad)\bm{u}}(\bm{k})=\sum_{\bm{j}+\bm{l}=\bm{k}}
(\hat{\bm{u}}(\bm{j})\cdot i\bm{l})\hat{\bm{u}}(\bm{l})$.
We define four BSDT channels:
\begin{enumerate}
\item \textbf{Enstrophy ($\delta_C^{\mathrm{NS}}$):}
  $\mathcal{E}(t)=\frac{1}{2}\norm{\bm{\omega}(t)}_{L^2}^2$
  where $\bm{\omega}=\grad\times\bm{u}$. This is the fluid
  analogue of credit camouflage: enstrophy measures deviation from
  laminar flow as Mahalanobis distance measures deviation from
  financial normality.

\item \textbf{Spectral cascade ($\delta_G^{\mathrm{NS}}$):}
  $E(k,t)=\sum_{|\bm{k}'|=k}|\hat{\bm{u}}(\bm{k}',t)|^2$. Energy
  transfer from low to high wavenumbers corresponds to feature-gap
  detection: motion into historically unoccupied spectral regions.

\item \textbf{Vorticity--strain alignment
  ($\delta_A^{\mathrm{NS}}$):}
  $\inner{\bm{\omega}}{S\bm{\omega}}/
  (\norm{\bm{\omega}}\norm{S\bm{\omega}})$
  where $S_{ij}=\frac{1}{2}(\partial_i u_j+\partial_j u_i)$.
  Anomalous alignment between vorticity and the middle eigenvector
  of the strain tensor is the turbulence analogue of activity
  anomaly \cite{ashurst1987,tsinober1992}.

\item \textbf{Temporal acceleration ($\delta_T^{\mathrm{NS}}$):}
  $\norm{\partial_t\bm{u}}_{L^2}^2$. Rapid temporal change in
  the velocity field corresponds to temporal novelty in the
  financial framework.
\end{enumerate}

The fluid blind-spot energy is
\begin{equation}\label{eq:ebs_ns}
  \Ebs^{\mathrm{NS}}(t) = w_1\mathcal{E}(t)
  + w_2\,\delta_G^{\mathrm{NS}}(t)
  + w_3\,\delta_A^{\mathrm{NS}}(t)
  + w_4\,\delta_T^{\mathrm{NS}}(t),
\end{equation}
with weights $w_i$ normalised so that
$\Ebs^{\mathrm{NS}}=1$ in the Stokes regime.

\subsection{Regularity Claim}\label{sec:ns:regularity}

\begin{theorem}[Conditional Regularity under Adaptive
Viscosity]\label{thm:ns_regularity}
Consider the modified Navier--Stokes system with adaptive viscosity
\begin{equation}\label{eq:ns_adaptive}
  \partial_t\bm{u}+(\bm{u}\cdot\grad)\bm{u}
  = -\grad p + \nu_{\mathrm{eff}}(t)\,\Delta\bm{u},\qquad
  \nu_{\mathrm{eff}}(t) = \nu_0\bigl(1+\gamma^*(\Ebs^{\mathrm{NS}}(t))\bigr),
\end{equation}
where $\gamma^*(E)=E/(E+\theta)$ is the optimal adaptive friction.
If the initial data $\bm{u}_0\in H^1(\T^3)$ and the enstrophy
satisfies
\begin{equation}\label{eq:enstrophy_bound}
  \sup_{t\in[0,T]}\mathcal{E}(t)\le
  C\nu_0^{-1}\norm{\bm{u}_0}_{H^1}^2,
\end{equation}
then the solution remains smooth on $[0,T]$:
$\bm{u}\in C^\infty(\T^3\times(0,T])$.
\end{theorem}

\begin{proof}[Proof sketch]
The enstrophy equation for the adaptive system is:
\begin{align}
  \frac{d\mathcal{E}}{dt}
  &= \underbrace{\int_{\T^3}\omega_i S_{ij}\omega_j\,dx}_{P
     \;\text{(production)}}
  - \underbrace{\nu_{\mathrm{eff}}(t)
    \norm{\grad\bm{\omega}}_{L^2}^2}_{D\;\text{(dissipation)}}.
    \label{eq:enstrophy_evolution}
\end{align}
When $\Ebs^{\mathrm{NS}}$ is large (indicating approach to
singularity), $\gamma^*(\Ebs^{\mathrm{NS}})\to 1$, so
$\nu_{\mathrm{eff}}\to 2\nu_0$. The enhanced dissipation competes
with the production term. By the Brezis--Gallouet inequality
\cite{bg1980}:
\[
  P \le C\norm{\bm{\omega}}_{L^2}^2\norm{\grad\bm{\omega}}_{L^2}
  \log^{1/2}(e+\norm{\Delta\bm{\omega}}_{L^2}/\norm{\bm{\omega}}_{L^2}).
\]
Under the enstrophy bound \eqref{eq:enstrophy_bound}, the BKM
criterion \cite{bkm1984}
$\int_0^T\norm{\bm{\omega}}_{L^\infty}\,dt<\infty$ is satisfied,
and classical regularity theory \cite{lady1969} extends the
solution.
\end{proof}

\begin{remark}
Theorem~\ref{thm:ns_regularity} does \emph{not} claim to resolve
the Millennium Prize problem. It establishes conditional regularity
under a \emph{modified} equation (adaptive viscosity), not the
original constant-viscosity Navier--Stokes. The contribution is the
structural connection: BSDT detects incipient singularity formation,
and adaptive damping provides a mechanism to suppress it.
\end{remark}

\subsection{The Production/Dissipation Halving
Principle}\label{sec:ns:pd}

The key numerical finding is that the adaptive viscosity maintains
the production/dissipation ratio $P/D$ at a universal value across
Reynolds numbers.

\begin{table}[htbp]
\centering
\caption{Production/dissipation ratio convergence under adaptive
viscosity. Pseudo-spectral Taylor--Green solver, $N=32$,
$T=1.0$, $\frac{2}{3}$-dealiasing.}
\label{tab:pd}
\begin{tabular}{@{}rcccc@{}}
\toprule
$\mathrm{Re}$ & $\nu_0$ & $P/D$ (const.\ $\nu$) &
$P/D$ (adaptive) & Ratio\\
\midrule
314   & 0.020 & 0.866 & 0.426 & 0.492\\
628   & 0.010 & 0.914 & 0.457 & 0.500\\
1257  & 0.005 & 0.938 & 0.474 & 0.505\\
3142  & 0.002 & 0.956 & 0.477 & 0.499\\
6283  & 0.001 & 0.971 & 0.485 & 0.500\\
\bottomrule
\end{tabular}
\end{table}

The ratio converges to $0.500\pm 0.005$ across two decades of
Reynolds number. This \emph{P/D halving principle} means the
adaptive viscosity exactly balances production against enhanced
dissipation, maintaining the system at the threshold between
enstrophy growth and decay. The mechanism is the fluid analogue of
the critical-manifold result (Theorem~\ref{thm:phase}(iii)):
adaptive friction holds the system at marginal stability.

\subsection{Depletion of Nonlinearity}\label{sec:ns:depletion}

The adaptive-viscosity mechanism effectively implements a
\emph{depletion of nonlinearity}: when enstrophy rises (signalling
approach to singularity), $\nu_{\mathrm{eff}}$ increases, which
increases dissipation, which reduces enstrophy, which reduces
$\nu_{\mathrm{eff}}$ back toward $\nu_0$. This feedback loop
implements the Lyapunov descent property (Theorem~C) in the setting
of~\eqref{eq:ns}.

The connection to classical regularity theory runs through the
Beale--Kato--Majda criterion \cite{bkm1984}: a smooth solution can
develop a singularity at time $T^*$ only if
$\int_0^{T^*}\norm{\bm{\omega}(\cdot,t)}_{L^\infty}\,dt=\infty$.
The adaptive viscosity bounds $\mathcal{E}(t)$ (and hence
$\norm{\bm{\omega}}_{L^2}$, which controls
$\norm{\bm{\omega}}_{L^\infty}$ via Sobolev embedding on $\T^3$),
preventing the integral from diverging.

\subsection{Connection to the Millennium Prize
Problem}\label{sec:ns:millennium}

The Millennium Prize \cite{fefferman2006} asks for a proof that
smooth solutions to \eqref{eq:ns} with $\nu>0$ exist globally, or
a demonstration of finite-time blowup. Our framework suggests three
strategies:
\begin{enumerate}
\item \textbf{A priori bounds on $\Ebs^{\mathrm{NS}}$.} If one
  can show that $\Ebs^{\mathrm{NS}}(t)\le C(\bm{u}_0,\nu)$ for
  all $t>0$ under constant viscosity, Theorem~\ref{thm:ns_regularity}
  implies global regularity. This would require proving that the
  natural Navier--Stokes dynamics are ``self-regulating'' in the
  BSDT sense.

\item \textbf{Vanishing adaptive correction.} If the adaptive
  correction $\gamma^*(\Ebs^{\mathrm{NS}})\to 0$ as $t\to\infty$
  (the fluid ``relaxes'' to a state where no extra viscosity is
  needed), the adaptive and constant-viscosity solutions converge,
  and regularity of the former implies regularity of the latter.

\item \textbf{Structural singularity obstruction.} The P/D halving
  principle (Table~\ref{tab:pd}) suggests that singularity
  formation requires $P/D\to\infty$---unbounded enstrophy
  production relative to dissipation. If the nonlinear term has an
  intrinsic tendency to deplete (as observed in numerical
  simulations of the full equations \cite{taylor1937}), this ratio
  may be bounded \emph{ab initio}, preventing blowup.
\end{enumerate}

These pathways are speculative. We emphasise that the current
results are \emph{conditional} (regularity under adaptive viscosity)
and \emph{numerical} (P/D halving observed computationally).
Converting them to unconditional \emph{a priori} bounds is an
open problem.


%% ====================================================================
%%  9.  DISCUSSION AND CONCLUSION
%% ====================================================================
\section{Discussion}\label{sec:discussion}

\subsection{Summary of Findings}

The paper establishes eight empirical findings within a unified
mathematical framework:

\begin{enumerate}
\item The U.S.\ banking system spends 68.6\% of time above the
  critical manifold $\Cman$---structurally over-coupled.

\item The MFLS signal leads the GFC by 6~quarters with
  AUROC $\ge 0.70$ on genuinely out-of-sample data.

\item All five scoring variants universally fail Granger
  causality---the expected signature of a phase-transition detector.

\item The Signed-LR variant discovers herding: negative
  $\beta_{\delta_T}=-1.38$ reveals that crises are preceded by
  \emph{decreased} temporal novelty.

\item Institution-level replication ($N=30$ banks) strengthens all
  results (AUROC = 0.805).

\item The Terra/Luna collapse ($r=0.996$) and ERCOT grid failure
  ($+72$\,h lead) validate cross-domain universality.

\item The framework transfers to Navier--Stokes: adaptive viscosity
  suppresses enstrophy growth with the P/D ratio converging to 0.500.

\item The welfare calibration produces GFC losses of 19.57~pp,
  consistent with the macro consensus.
\end{enumerate}

\subsection{Limitations}

\begin{enumerate}
\item \textbf{Pseudo-spectral resolution.} The NS simulations use
  $N=32$ modes. Higher-resolution simulations ($N\ge 256$) are
  needed to confirm P/D halving at fully turbulent Reynolds numbers.

\item \textbf{Conditional regularity.} The NS result is for the
  \emph{modified} equation (adaptive viscosity), not the original
  constant-viscosity problem.

\item \textbf{Gradient anti-alignment.} The empirical anti-alignment
  ($\costheta=-0.335$) on real data shows that $\grad\Ebs$ and
  $\grad\Phi$ measure complementary rather than parallel information.
  The equivalence holds in the norm sense (Theorem~C) but the
  directions differ.

\item \textbf{Data frequency.} FDIC call reports are quarterly;
  higher-frequency data (daily balance sheets) would sharpen lead
  times.

\item \textbf{False alarm rate.} The 35\% FAR reflects persistent
  super-criticality rather than model noise, but reduces operational
  precision.

\item \textbf{Small $N$.} The primary panel has $N=7$ institution
  types. The $N=30$ extension partially addresses this.
\end{enumerate}

\subsection{Open Problems}

\begin{enumerate}
\item Establish unconditional \emph{a priori} bounds on
  $\Ebs^{\mathrm{NS}}$ under constant viscosity.

\item Prove the P/D halving principle analytically.

\item Determine whether the separation condition
  (Assumption~\ref{ass:sep}) is necessary or merely sufficient for
  detection-based stabilisation.

\item Extend the OCO regret bound to the nonlinear-cost setting
  $c_t(\gamma)=\Phi(X^{t+1}(\gamma))-\Phi(X^t)$ with exact
  convexity (not first-order approximation).

\item Connect the Pigouvian interpretation to optimal currency-area
  theory: is the CCyB path optimal across heterogeneous
  jurisdictions?

\item Extend the ERCOT counterfactual to a full N-1 contingency
  analysis with realistic grid topology.

\item Characterise the basin of attraction of $\Xst$ under
  adaptive friction as a function of the noise intensity $\sigma$.
\end{enumerate}

\subsection{Policy Implications}

For \emph{central banks}: the framework provides a theoretically
grounded alternative to the credit-to-GDP gap for CCyB calibration,
with explicit welfare costs and real-time implementability.
The negative $\delta_A$ finding implies that low-volatility periods
should trigger \emph{increased}, not decreased, buffers.

For the \emph{Navier--Stokes community}: the transfer demonstrates
that modern financial mathematics---online learning, adaptive
control, Lyapunov analysis---may contribute tools for classical
PDE regularity questions. The P/D halving principle, if confirmed
at higher resolution, provides a concrete target for analytical
proof.


%% ====================================================================
%%  APPENDIX
%% ====================================================================
\begin{appendix}

\section{Full Proofs}\label{app:proofs}

\begin{proof}[Proof of Lemma~\ref{lem:upper} (full)]
From Lemma~\ref{lem:grad_ebs} and the triangle inequality:
\begin{align*}
  \norm{\grad\Ebs(X)}
  &\le \sum_{i=1}^{4}|\psi_i'(\delta_i)|\norm{\grad_X\delta_i}
  + \norm{F_{\mathrm{pair}}}\\
  &\le \sum_{i=1}^{4}K_i\norm{X{-}\Xst}+\Lpair\norm{X{-}\Xst}\\
  &\le (\textstyle\sum_i K_i+\Lpair)\norm{X{-}\Xst}.
\end{align*}
Since $\norm{\grad\Phi(X)}\ge\lmin(H^*)^{1/2}\norm{X{-}\Xst}$
near $\Xst$, we obtain $\norm{\grad\Ebs}\le c_2\norm{\grad\Phi}$
with $c_2=(\sum_i K_i+\Lpair)/\lmin(H^*)^{1/2}$.
\end{proof}

\begin{proof}[Proof of Lemma~\ref{lem:lower} (full)]
Choose $v=\grad\Phi(X)/\norm{\grad\Phi(X)}$. By
Assumption~\ref{ass:sep} and the chain rule:
\[
  \inner{\grad\Ebs^{\mathrm{dev}}}{v}
  = \sum_i\psi_i'(\delta_i)\inner{\grad_X\delta_i}{v}
  \ge \sum_i\psi_i'(\delta_i)|\inner{\grad_X\delta_i}{v}|.
\]
By Assumption~\ref{ass:sep},
$|\inner{\grad_X\delta_i}{v}|\ge\nu\norm{\grad\Phi}$ for some $i$.
Since $\psi_i'\ge p_{\min}>0$ for $\delta_i>0$:
$\norm{\grad\Ebs}\ge\inner{\grad\Ebs}{v}\ge
p_{\min}\nu\norm{\grad\Phi}$.
Set $c_1=p_{\min}\nu$.
\end{proof}

\begin{proof}[Proof of Theorem~\ref{thm:descent} (full)]
Define $V(t)=\Phi(X(t))$. Differentiating:
$\dot{V}=\grad\Phi^\top\dot{X}=\grad\Phi^\top(-\grad\Phi)
=-\norm{\grad\Phi}^2\le 0$. Since $V(t)$ is non-increasing and
bounded below by $\Phi(\Xst)$, it converges to $V_\infty$. By
LaSalle \cite{lasalle1960}, every $\omega$-limit point $X_\infty$
satisfies $\grad\Phi(X_\infty)=0$.
\end{proof}

\begin{proof}[Proof of Theorem~\ref{thm:las} (full)]
Write $z(t)=X(t)-\Xst$. Then $\dot{z}=-H^*z+O(\norm{z}^2)$ and
$\frac{d}{dt}\norm{z}^2\le-2(\alpha-\Lpair)\norm{z}^2$ by Gronwall,
giving $\norm{z(t)}\le\norm{z(0)}e^{-(\alpha-\Lpair)t}$.
\end{proof}

\begin{proof}[Proof of Theorem~\ref{thm:phase} (full)]
\textit{Part~(ii).} Above $\Cman$, let $u$ be the unit eigenvector
of $D^2\Phi_{\mathrm{pair}}(X^+)$ for $\lambda_u>\alpha$. Define
$W(t)=\frac{1}{2}\inner{X(t){-}\Xst}{u}^2$. Then:
$\dot{W}=(-\alpha+\lambda_u)W+O(\norm{z}^3)=cW+O(\norm{z}^3)$
with $c=\lambda_u-\alpha>0$. By forward Gronwall:
$W(t)\ge W(0)e^{ct/2}$, growing exponentially. No fixed
$\bar{\gamma}$ prevents this.
\end{proof}

\begin{proof}[Proof of Proposition~\ref{prop:regret} (full)]
Standard OGD \cite{hazan2016oco} with step
$\eta_\gamma=(\gamma_{\max}{-}\gamma_{\min})/\sqrt{T}$:
\begin{align*}
  R_T &\le \frac{(\gamma_{\max}{-}\gamma_{\min})^2}{2\eta_\gamma}
  + \frac{\eta_\gamma}{2}\sum_t\norm{\grad_\gamma c_t}^2\\
  &\le L_c(\gamma_{\max}{-}\gamma_{\min})\sqrt{T}.
\end{align*}
\end{proof}

\begin{proof}[Proof of Theorem~\ref{thm:optimal} (full)]
First-order condition of the Bellman equation w.r.t.\ $\gamma$:
$2\kappa\gamma=\beta\partial_\gamma\E[V(X')|X,\gamma]$. Using
$\partial_\gamma X'=-\eta\grad E(X)$ and
$\partial_{X'}V\approx\grad E(X')/(1{-}\beta)$:
$2\kappa\gamma^*\approx\beta\eta\norm{\grad E}^2/(1{-}\beta)$,
giving $\gamma^*\approx\beta\eta\norm{\grad E}^2/[2\kappa(1{-}\beta)]$.
The saturation factor $E/(E{+}\theta)$ ensures $\gamma^*\to 0$ as
$E\to 0$.
\end{proof}

\section{Computational Details}\label{app:compute}

\subsection{GravityEngine Parameters}

\begin{table}[htbp]
\centering
\caption{GravityEngine parameter table.}
\label{tab:params}
\begin{tabular}{llll}
\toprule
Parameter & Symbol & Default & Role\\
\midrule
Radial spring     & $\alpha$   & 0.1     & Mean-reversion\\
Pairwise coupling & $\gamma$   & 1.0     & Interaction scale\\
Gaussian range    & $\sigma$   & 1.0     & Attraction length\\
Repulsion coeff.  & $\lambda$  & 0.1     & Prevents collapse\\
Force clamp       & $F_{\max}$ & 100     & Bounds Jacobian\\
Step size         & $\eta$     & 0.02    & Euler step\\
\bottomrule
\end{tabular}
\end{table}

\subsection{Navier--Stokes Solver}

The pseudo-spectral solver uses $N^3$ Fourier collocation points
on $[0,2\pi]^3$ with $\frac{2}{3}$-dealiasing. Time integration is
semi-implicit: exact viscous integration (integrating factor) with
explicit nonlinear term. The adaptive viscosity
$\nu_{\mathrm{eff}}=\nu_0(1+\gamma^*(\Ebs^{\mathrm{NS}}))$ is
computed at each diagnostic step. Each time step costs
$O(N^3\log N)$ for FFTs. Extended simulation at $\mathrm{Re}=1257$
($\nu_0=0.005$, $N=32$, $T=2.0$):

\begin{table}[htbp]
\centering
\caption{Extended NS simulation at $\mathrm{Re}=1257$.}
\label{tab:ns_extended}
\begin{tabular}{@{}lcc@{}}
\toprule
Quantity & Const.\ $\nu$ & Adaptive $\nu(\Ebs)$\\
\midrule
Peak enstrophy     & 0.5106 & 0.4625\\
Peak $\norm{\bm{\omega}}_\infty$ & 2.000 & 2.000\\
BKM integral       & 3.304  & 3.235\\
Energy decay (\%)  & 6.56   & 12.47\\
$\nu_{\mathrm{eff}}$ range & $[0.005,0.005]$
  & $[0.005,0.010]$\\
$\gamma^*_{\max}$  & ---    & 0.995\\
\bottomrule
\end{tabular}
\end{table}

The adaptive mechanism reduces peak enstrophy by 9.4\% and nearly
doubles the energy dissipation rate (12.47\% vs.\ 6.56\%) while
maintaining smoothness.

\section{Data Sources}\label{app:data}

All banking features are drawn from the FDIC SDI quarterly call
reports via the public REST API at
\texttt{https://banks.data.fdic.gov/api/financials}. Feature $f_5$
(yield-curve slope) is the FRED T10Y2Y series. The normal reference
period (1994-Q1 to 2003-Q4) is chosen to be post-S\&L resolution,
pre-housing boom, and long enough for stable covariance estimation
($T_{\mathrm{ref}}=40\gg d=6$). The full replication package is
available at
\texttt{https://github.com/segetii/fintech}.

\end{appendix}


%% ====================================================================
%%  BIBLIOGRAPHY
%% ====================================================================
\bibliographystyle{siam}
\begin{thebibliography}{99}

\bibitem{acemoglu2015systemic}
D.~Acemoglu, A.~Ozdaglar, and A.~Tahbaz-Salehi,
\emph{Systemic risk and stability in financial networks},
Amer.\ Econ.\ Rev., 105 (2015), pp.~564--608.

\bibitem{adrian2016covar}
T.~Adrian and M.~K.~Brunnermeier,
\emph{CoVaR},
Amer.\ Econ.\ Rev., 106 (2016), pp.~1705--1741.

\bibitem{ashurst1987}
W.~T.~Ashurst, A.~R.~Kerstein, R.~M.~Kerr, and C.~H.~Gibson,
\emph{Alignment of vorticity and scalar gradient with strain rate
in simulated Navier--Stokes turbulence},
Phys.\ Fluids, 30 (1987), pp.~2343--2353.

\bibitem{bkm1984}
J.~T.~Beale, T.~Kato, and A.~Majda,
\emph{Remarks on the breakdown of smooth solutions for the 3-D
Euler equations},
Comm.\ Math.\ Phys., 94 (1984), pp.~61--66.

\bibitem{bertsekas2012dynamic}
D.~P.~Bertsekas,
\emph{Dynamic Programming and Optimal Control},
4th ed., Athena Scientific, Belmont, MA, 2012.

\bibitem{bg1980}
H.~Brezis and T.~Gallouet,
\emph{Nonlinear Schr\"odinger evolution equations},
Nonlinear Anal., 4 (1980), pp.~677--681.

\bibitem{brownlees2017srisk}
C.~Brownlees and R.~F.~Engle,
\emph{SRISK: A conditional capital shortfall measure of systemic
risk},
Rev.\ Financ.\ Stud., 30 (2017), pp.~48--79.

\bibitem{brunnermeier2009deciphering}
M.~K.~Brunnermeier,
\emph{Deciphering the liquidity and credit crunch 2007--2008},
J.\ Econ.\ Perspectives, 23 (2009), pp.~77--100.

\bibitem{elliott2014financial}
M.~Elliott, B.~Golub, and M.~O.~Jackson,
\emph{Financial networks and contagion},
Amer.\ Econ.\ Rev., 104 (2014), pp.~3115--3153.

\bibitem{fefferman2006}
C.~L.~Fefferman,
\emph{Existence and smoothness of the Navier--Stokes equation},
Clay Mathematics Institute Millennium Prize Problems, 2000.

\bibitem{greenwald2020financial}
D.~L.~Greenwald, \`O.~Jord\`a, M.~Schularick, and
A.~M.~Taylor,
\emph{Financial crisis recoveries},
AEA Papers and Proc., 110 (2020), pp.~124--129.

\bibitem{hazan2016oco}
E.~Hazan,
\emph{Introduction to Online Convex Optimization},
Found.\ Trends Optim., 2 (2016), pp.~157--325.

\bibitem{khalil2002nonlinear}
H.~K.~Khalil,
\emph{Nonlinear Systems},
3rd ed., Prentice Hall, Upper Saddle River, NJ, 2002.

\bibitem{km2006}
C.~E.~Kenig and F.~Merle,
\emph{Global well-posedness, scattering and blow-up for the
energy-critical, focusing, non-linear Schr\"odinger equation in
the radial case},
Invent.\ Math., 166 (2006), pp.~645--675.

\bibitem{lady1969}
O.~A.~Ladyzhenskaya,
\emph{The Mathematical Theory of Viscous Incompressible Flow},
2nd ed., Gordon and Breach, New York, 1969.

\bibitem{lasalle1960}
J.~P.~LaSalle,
\emph{Some extensions of Liapunov's second method},
IRE Trans.\ Circuit Theory, 7 (1960), pp.~520--527.

\bibitem{ledoit2004}
O.~Ledoit and M.~Wolf,
\emph{A well-conditioned estimator for large-dimensional covariance
matrices},
J.\ Multivariate Anal., 88 (2004), pp.~365--411.

\bibitem{lohmiller1998contraction}
W.~Lohmiller and J.-J.~E.~Slotine,
\emph{On contraction analysis for non-linear systems},
Automatica, 34 (1998), pp.~683--696.

\bibitem{may2008complex}
R.~M.~May, S.~A.~Levin, and G.~Sugihara,
\emph{Complex systems: Ecology for bankers},
Nature, 451 (2008), pp.~893--895.

\bibitem{odeyemi2026amttp}
O.~Odeyemi,
\emph{{AMTTP}: A four-layer architecture for deterministic
compliance enforcement in institutional {DeFi}},
TechRxiv, 2026.

\bibitem{odeyemi2026bsdt}
O.~Odeyemi,
\emph{Blind spot decomposition theory: A mathematical framework for
characterising and correcting missed fraud in machine learning
systems},
Zenodo, 2026.

\bibitem{shalev2012online}
S.~Shalev-Shwartz,
\emph{Online learning and online convex optimization},
Found.\ Trends Mach.\ Learn., 4 (2012), pp.~107--194.

\bibitem{taylor1937}
G.~I.~Taylor and A.~E.~Green,
\emph{Mechanism of the production of small eddies from large ones},
Proc.\ R.\ Soc.\ Lond.\ A, 158 (1937), pp.~499--521.

\bibitem{trefethen1997numerical}
L.~N.~Trefethen and D.~Bau,
\emph{Numerical Linear Algebra},
SIAM, Philadelphia, 1997.

\bibitem{tsinober1992}
A.~Tsinober, E.~Kit, and T.~Dracos,
\emph{Experimental investigation of the field of velocity gradients
in turbulent flows},
J.\ Fluid Mech., 242 (1992), pp.~169--192.

\end{thebibliography}

\end{document}
